
\chapter{アルゴリズム的基礎}\label{chap:algorithmic}

本章では，離散的な Kantorovich 問題を\textbf{厳密に}解くための
アルゴリズムを扱う．
主問題（線形計画）と双対問題の構造を活用する3つのアプローチ
--- ネットワーク単体法，双対上昇法，オークションアルゴリズム ---
を紹介する．

これらのアルゴリズムの起源は第二次世界大戦前後に遡る．
Dantzig~\cite{Dantzig1951} が単体法を開発した主な動機の一つが
まさに輸送問題であり，最適輸送は線形計画法の発展と密接に関わってきた．


% ============================================================
\section{Kantorovich 問題の線形計画としての表現}\label{sec:kant-lp}

第\ref{chap:theoretical}章で導入した Kantorovich 問題
\[
  \mathbf{L_C}(\mathbf{a}, \mathbf{b})
  = \min_{\mathbf{P} \in \mathbf{U}(\mathbf{a}, \mathbf{b})}
  \inner{\mathbf{C}}{\mathbf{P}}
\]
を標準的な線形計画の形に書き直そう．

行列 $\mathbf{P} \in \R^{n \times m}$ を
列ごとに縦に並べたベクトル $\mathbf{p} \in \R^{nm}$ に変換し，
制約行列
\[
  \mathbf{A} =
  \begin{bmatrix}
    \mathbb{1}_n^\top \otimes \mathbb{I}_m \\
    \mathbb{I}_n \otimes \mathbb{1}_m^\top
  \end{bmatrix}
  \in \R^{(n+m) \times nm}
\]
を用いると，Kantorovich 問題は標準形の線形計画
\begin{equation}\label{eq:lp-standard}
  \min_{\substack{\mathbf{p} \in \R^{nm}_+ \\
  \mathbf{A}\mathbf{p} = \begin{bsmallmatrix} \mathbf{a} \\ \mathbf{b} \end{bsmallmatrix}}}
  \mathbf{c}^\top \mathbf{p}
\end{equation}
と書ける．
ここで $\mathbf{c} \in \R^{nm}$ は
コスト行列 $\mathbf{C}$ を列ごとに並べたベクトルである．

対応する双対問題は
$\mathbf{h} = \bigl(\begin{smallmatrix} \mathbf{f} \\ \mathbf{g} \end{smallmatrix}\bigr) \in \R^{n+m}$ として
\begin{equation}\label{eq:lp-dual}
  \max_{\substack{\mathbf{h} \in \R^{n+m} \\
  \mathbf{A}^\top \mathbf{h} \leq \mathbf{c}}}
  \begin{bmatrix} \mathbf{a} \\ \mathbf{b} \end{bmatrix}^\top \mathbf{h}
\end{equation}
であり，第\ref{chap:theoretical}章の
Kantorovich 双対（命題\ref{prop:kant-dual}）に一致する．

\begin{remark}{制約の冗長性}{constraint-redundancy}
行の和制約 $n$ 本と列の和制約 $m$ 本の計 $n+m$ 本のうち，
1本は他の制約から導かれる
（$\sum_i \mathbf{a}_i = \sum_j \mathbf{b}_j = 1$ であるため）．
よって実質的な等式制約は $n+m-1$ 本である．
この冗長性は，双対変数 $(\mathbf{f}, \mathbf{g})$ に
定数の自由度が残ることに対応する．
\end{remark}


% ============================================================
\section{$\mathbf{C}$変換}\label{sec:c-transform}

双対問題の構造を活かして，
変数の一方を固定したときに
もう一方を最適に選ぶ操作を\textbf{$\mathbf{C}$変換}と呼ぶ．
これは双対問題の解法やオークションアルゴリズムの
基本的な構成要素となる．

\begin{definition}{$\mathbf{C}$変換と $\bar{\mathbf{C}}$変換}{c-transform}
双対ポテンシャル $\mathbf{f} \in \R^n$ に対し，
その\textbf{$\mathbf{C}$変換} $\mathbf{f}^{\mathbf{C}} \in \R^m$ を
\begin{equation}\label{eq:c-transform}
  (\mathbf{f}^{\mathbf{C}})_j
  \defeq \min_{i \in \{1,\ldots,n\}} \mathbf{C}_{i,j} - \mathbf{f}_i
\end{equation}
で定義する．
同様に，$\mathbf{g} \in \R^m$ の
$\bar{\mathbf{C}}$\textbf{変換} $\mathbf{g}^{\bar{\mathbf{C}}} \in \R^n$ を
\begin{equation}\label{eq:cbar-transform}
  (\mathbf{g}^{\bar{\mathbf{C}}})_i
  \defeq \min_{j \in \{1,\ldots,m\}} \mathbf{C}_{i,j} - \mathbf{g}_j
\end{equation}
で定義する．
\end{definition}

$\mathbf{C}$変換の意味は明快である：
$\mathbf{f}$ を固定したとき，双対制約
$\mathbf{f}_i + \mathbf{g}_j \leq \mathbf{C}_{i,j}$
をすべての $i$ について満たす $\mathbf{g}_j$ の最大値が
$(\mathbf{f}^{\mathbf{C}})_j$ である．
したがって，$\mathbf{f}$ に対する最良の相方として
$\mathbf{g} = \mathbf{f}^{\mathbf{C}}$ を選ぶことで
双対目的関数が改善される：
\[
  \inner{\mathbf{f}}{\mathbf{a}} + \inner{\mathbf{g}}{\mathbf{b}}
  \leq
  \inner{\mathbf{f}}{\mathbf{a}} + \inner{\mathbf{f}^{\mathbf{C}}}{\mathbf{b}}.
\]

$\mathbf{C}$変換と $\bar{\mathbf{C}}$変換を交互に適用すると
双対目的関数は単調非減少になるが，有限回で不動点に到達する
（3回の変換で安定する）：
\[
  \mathbf{f}^{\mathbf{C}\bar{\mathbf{C}}\mathbf{C}} = \mathbf{f}^{\mathbf{C}}.
\]

\begin{example}{$\mathbf{C}$変換の計算}{c-transform-example}
$n = m = 2$，
$\mathbf{C} = \bigl(\begin{smallmatrix} 1 & 3 \\ 4 & 2 \end{smallmatrix}\bigr)$，
$\mathbf{f} = (0,\; 0)$ とする．
\begin{align*}
  (\mathbf{f}^{\mathbf{C}})_1
  &= \min(1-0,\; 4-0) = 1, \\
  (\mathbf{f}^{\mathbf{C}})_2
  &= \min(3-0,\; 2-0) = 2.
\end{align*}
よって $\mathbf{f}^{\mathbf{C}} = (1, 2)$ であり，
制約 $\mathbf{f}_i + \mathbf{g}_j \leq \mathbf{C}_{i,j}$
はすべて満たされる（等号が成立するのは
$(1,1)$: $0+1=1$ と $(2,2)$: $0+2=2$）．
\end{example}


% ============================================================
\section{相補性条件と最適性判定}\label{sec:complementary}

第\ref{chap:theoretical}章で述べた相補性条件を
アルゴリズムの視点から整理する．

\begin{definition}{相補性}{complementarity-def}
行列 $\mathbf{P} \in \R^{n \times m}$ とベクトル対
$(\mathbf{f}, \mathbf{g})$ がコスト行列 $\mathbf{C}$ に関して
\textbf{相補的} (complementary) であるとは，
\[
  \mathbf{P}_{i,j} > 0
  \quad \Longrightarrow \quad
  \mathbf{f}_i + \mathbf{g}_j = \mathbf{C}_{i,j}
\]
がすべての $(i,j)$ について成り立つことをいう．
\end{definition}

\begin{proposition}{相補的な実行可能解は最適}{comp-optimal}
$\mathbf{P} \in \mathbf{U}(\mathbf{a}, \mathbf{b})$ と
$(\mathbf{f}, \mathbf{g}) \in \mathbf{R}(\mathbf{C})$ が
相補的ならば，$\mathbf{P}$ は主問題の最適解，
$(\mathbf{f}, \mathbf{g})$ は双対問題の最適解である．
\end{proposition}

\begin{proof}
弱双対性より
$\mathbf{L_C}(\mathbf{a}, \mathbf{b})
\leq \inner{\mathbf{P}}{\mathbf{C}}$ かつ
$\inner{\mathbf{f}}{\mathbf{a}} + \inner{\mathbf{g}}{\mathbf{b}}
\leq \mathbf{L_C}(\mathbf{a}, \mathbf{b})$．
相補性から
\[
  \inner{\mathbf{P}}{\mathbf{C}}
  = \inner{\mathbf{P}}{\mathbf{f} \oplus \mathbf{g}}
  = \inner{\mathbf{f}}{\mathbf{a}} + \inner{\mathbf{g}}{\mathbf{b}}
\]
が成り立つので，不等式が連鎖的に等号になり両方が最適である．
\end{proof}

この命題は，以下で述べるすべてのアルゴリズムの
正当性の基盤となる：
\textbf{実行可能かつ相補的な主・双対の組を見つければ，
それが最適解}である．


% ============================================================
\section{輸送多面体の頂点構造}\label{sec:transport-polytope}

カップリングの集合 $\mathbf{U}(\mathbf{a}, \mathbf{b})$ は
$\R^{n \times m}$ の凸多面体（\textbf{輸送多面体}）である．
線形計画の基本定理から，最適解はこの多面体の
\textbf{頂点}（端点）で達成される．
輸送多面体の頂点は，次の意味でスパースである．

\begin{proposition}{頂点の木構造}{extremal-tree}
$\mathbf{P}$ が $\mathbf{U}(\mathbf{a}, \mathbf{b})$ の
頂点ならば，$\mathbf{P}_{i,j} > 0$ となるインデックスの集合
$S(\mathbf{P}) = \{(i,j) : \mathbf{P}_{i,j} > 0\}$ を辺集合とする
二部グラフ $G(\mathbf{P})$ は閉路を含まない（森である）．
特に，非零成分の個数は高々 $n + m - 1$ 個である．
\end{proposition}

\begin{proof}
背理法による．$G(\mathbf{P})$ が閉路を含むと仮定すると，
閉路に沿って交互に $+\varepsilon$, $-\varepsilon$ を加えた
2つのカップリング $\mathbf{Q}$, $\mathbf{R}$ が構成でき，
$\mathbf{P} = (\mathbf{Q} + \mathbf{R})/2$ となるが，
$\mathbf{Q} \neq \mathbf{R}$ であるから
$\mathbf{P}$ が頂点であることに矛盾する．
閉路のない $n + m$ 頂点のグラフの辺数は高々 $n + m - 1$ である．
\end{proof}

\begin{remark}{北西隅法}{nw-corner}
輸送多面体の頂点を簡便に一つ求めるヒューリスティックとして
\textbf{北西隅法} (North-West corner rule) がある．
$\mathbf{P}_{1,1} = \min(\mathbf{a}_1, \mathbf{b}_1)$ と設定し，
行または列の制約が飽和するごとにインデックスを進めて
$\mathbf{P}_{n,m}$ まで値を決定する．
得られる $\mathbf{P}$ は高々 $n + m - 1$ 個の非零成分を持ち，
$\mathbf{U}(\mathbf{a}, \mathbf{b})$ の頂点になる．

この解は最適とは限らないが，
ネットワーク単体法の初期解として利用できる．
\end{remark}

\begin{example}{北西隅法}{nw-example}
$\mathbf{a} = (0.2,\; 0.5,\; 0.3)$,
$\mathbf{b} = (0.5,\; 0.1,\; 0.4)$ とする．

\[
  \mathbf{P}^{\mathrm{NW}} =
  \begin{pmatrix}
    0.2 & 0 & 0 \\
    0.3 & 0.1 & 0.1 \\
    0 & 0 & 0.3
  \end{pmatrix}
\]

非零成分は 4 個であり，$n + m - 1 = 5$ 以下を満たす．
\end{example}


% ============================================================
\section{ネットワーク単体法}\label{sec:network-simplex}

\textbf{ネットワーク単体法}は，
輸送問題の木構造を活用した単体法の特殊版であり，
実用上最も広く使われる厳密解法の一つである．

\subsection{アルゴリズムの概要}\label{subsec:ns-overview}

アルゴリズムは以下のサイクルを繰り返す：

\begin{enumerate}
  \item \textbf{初期化}：
    北西隅法等で頂点解 $\mathbf{P}$ を求め，
    そのサポートグラフ $G(\mathbf{P})$（木）を構成する．
  \item \textbf{双対変数の計算}：
    $G(\mathbf{P})$ の辺 $(i,j)$ について
    $\mathbf{f}_i + \mathbf{g}_j = \mathbf{C}_{i,j}$
    を満たす $(\mathbf{f}, \mathbf{g})$ を求める．
    木の適当な根を選び $\mathbf{f}_{\mathrm{root}} = 0$ として，
    幅優先・深さ優先探索で順に値を決定する．
  \item \textbf{最適性の判定}：
    すべての $(i,j)$ について
    $\mathbf{f}_i + \mathbf{g}_j \leq \mathbf{C}_{i,j}$
    ならば最適であり，終了する．
  \item \textbf{ピボット操作}：
    $\mathbf{f}_i + \mathbf{g}_j > \mathbf{C}_{i,j}$ なる
    違反辺 $(i,j)$ を $G$ に追加すると閉路ができる．
    閉路に沿って流量を調整し，
    流量がゼロになった辺を除去して再び木に戻す．
\end{enumerate}

\begin{figure}[htbp]
\centering
\begin{tikzpicture}[>=Stealth,
  box/.style={draw, rounded corners, minimum width=5em, minimum height=2em,
              font=\small, align=center}]
  \node[box] (init) at (0, 0) {初期化\\(北西隅法)};
  \node[box] (dual) at (4, 0) {双対変数\\を計算};
  \node[box] (check) at (8, 0) {最適性\\の判定};
  \node[box] (pivot) at (8, -2) {ピボット\\操作};
  \node[font=\small] (end) at (12, 0) {最適解};
  \draw[->] (init) -- (dual);
  \draw[->] (dual) -- (check);
  \draw[->] (check) -- (end) node[midway, above]{\scriptsize Yes};
  \draw[->] (check) -- (pivot) node[midway, right]{\scriptsize No};
  \draw[->] (pivot) -| (dual);
\end{tikzpicture}
\caption{ネットワーク単体法の処理フロー．}
\label{fig:network-simplex-flow}
\end{figure}

\subsection{ピボット操作の詳細}\label{subsec:ns-pivot}

ステップ 4 の操作をもう少し詳しく説明する．

$\mathbf{f}_i + \mathbf{g}_j > \mathbf{C}_{i,j}$ を満たす
違反辺 $(i, j)$ を木 $G$ に追加すると，
$G$ にちょうど一つの閉路ができる．
この閉路上の辺を，$(i,j)$ を含む方向を「正の辺」，
逆方向を「負の辺」と分類する．

負の辺のうち流量が最小のもの
$\theta = \min\{\mathbf{P}_{i',j'} : (i',j') \text{ は負の辺}\}$
を見つけ，正の辺に $+\theta$，負の辺に $-\theta$ を加える．
流量がゼロになった辺を $G$ から除去すると再び木に戻る．

この操作により，主問題のコストは
\[
  \inner{\tilde{\mathbf{P}}}{\mathbf{C}} - \inner{\mathbf{P}}{\mathbf{C}}
  = \theta \bigl(\mathbf{C}_{i,j} - (\mathbf{f}_i + \mathbf{g}_j)\bigr)
  < 0
\]
と厳密に減少する（$\theta > 0$ のとき）．

\subsection{計算量}\label{subsec:ns-complexity}

ネットワーク単体法の多項式時間性は
Orlin~(1997) により初めて証明された．
その後の改良により，計算量は
\[
  O\bigl((n+m)\, nm\, \log(n+m)\, \log\bigl((n+m)\|\mathbf{C}\|_\infty\bigr)\bigr)
\]
に抑えられることが知られている．
実用的には，後述のエントロピー正則化法より遅いが，
\textbf{厳密解}が得られる利点がある．


% ============================================================
\section{双対上昇法}\label{sec:dual-ascent}

ネットワーク単体法が主問題の頂点解を反復的に改善するのに対し，
\textbf{双対上昇法}は
双対実行可能解 $(\mathbf{f}, \mathbf{g})$ を
反復的に改善して最適に到達する．

\subsection{基本的な考え方}\label{subsec:da-idea}

双対実行可能解 $(\mathbf{f}, \mathbf{g}) \in \mathbf{R}(\mathbf{C})$
に対し，$\mathbf{f}_i + \mathbf{g}_j = \mathbf{C}_{i,j}$
が等号で成立する辺の集合を\textbf{均衡辺} (balanced edge) と呼ぶ．
それ以外の辺（$\mathbf{f}_i + \mathbf{g}_j < \mathbf{C}_{i,j}$）を
\textbf{非活性辺}と呼ぶ．

双対上昇法では，
均衡辺のみを使って主制約を満たすカップリング $\mathbf{P}$
が構成できるかどうかを判定し，
構成できなければ $(\mathbf{f}, \mathbf{g})$ を修正して
双対目的関数を増加させる．

\subsection{最大流との関係}\label{subsec:da-maxflow}

具体的には，均衡辺のみからなる二部グラフに
仮想的な始点 $s$ と終点 $t$ を追加し，
$s$ から各 $i$ に容量 $\mathbf{a}_i$ の辺を，
各 $j$ から $t$ に容量 $\mathbf{b}_j$ の辺を張る．
この容量付きグラフ上の\textbf{最大流}が 1（= 全質量）であれば，
得られたフローがそのまま主問題の最適解 $\mathbf{P}$ を与える．

最大流が 1 未満の場合は，
Ford--Fulkerson のラベリングアルゴリズムにより
飽和していないノードの集合 $S \subset \{1,\ldots,n\}$，
$S' \subset \{1,\ldots,m\}$ が特定される．
$\mathbf{f}$ の $S$ 成分を $\varepsilon$ だけ増やし，
$\mathbf{g}$ の $S'$ 成分を $\varepsilon$ だけ減らす更新
\[
  (\tilde{\mathbf{f}}, \tilde{\mathbf{g}})
  = (\mathbf{f}, \mathbf{g})
  + \varepsilon(\mathbb{1}_S, -\mathbb{1}_{S'})
\]
を行う（$\varepsilon$ は双対実行可能性を保つ最大値）．
この更新により双対目的関数は厳密に増加する．

\subsection{Hungarian 法}\label{subsec:hungarian}

一様な割当問題
$\mathbf{a} = \mathbf{b} = \mathbb{1}_n / n$ の場合，
上記の双対上昇法は
\textbf{Hungarian 法}（Kuhn, 1955）に帰着される．
その計算量は $O(n^3)$ であり，
割当問題に対する最良の厳密解法の一つである．


% ============================================================
\section{オークションアルゴリズム}\label{sec:auction}

\textbf{オークションアルゴリズム}
（Bertsekas, 1981）は，
割当問題に特化した双対上昇法であり，
経済的なオークションの比喩で理解できる．

\subsection{問題設定と直感}\label{subsec:auction-intuition}

$n$ 人の入札者（$\alpha$ の点に対応）が
$n$ 個の品物（$\beta$ の点に対応）を競り落とすことを考える．
品物 $j$ には\textbf{価格} $\mathbf{g}_j$ が付いており，
入札者 $i$ にとって品物 $j$ の正味の利得は
$\mathbf{C}_{i,j} - \mathbf{g}_j$（小さいほど良い）である．

\subsection{アルゴリズムの手順}\label{subsec:auction-steps}

アルゴリズムは，
未割当の入札者集合 $S$，部分割当 $\xi$，
価格ベクトル $\mathbf{g}$ の三つ組を反復的に更新する．
各反復で，未割当の入札者 $i \notin S$ が以下を行う：

\begin{enumerate}
  \item \textbf{入札}：
    $i$ にとっての最良の品物 $j_i^1$ と次善の品物 $j_i^2$ を求める：
    \[
      j_i^1 \in \argmin_{j} \mathbf{C}_{i,j} - \mathbf{g}_j,
      \quad
      j_i^2 \in \argmin_{j \neq j_i^1} \mathbf{C}_{i,j} - \mathbf{g}_j.
    \]

  \item \textbf{価格更新}：
    品物 $j_i^1$ の価格を，最良と次善のコスト差に
    $\varepsilon$ を加えた分だけ引き下げる：
    \begin{equation}\label{eq:auction-price}
      \mathbf{g}_{j_i^1} \leftarrow
      \mathbf{C}_{i,j_i^1} - (\mathbf{C}_{i,j_i^2} - \mathbf{g}_{j_i^2})
      - \varepsilon.
    \end{equation}

  \item \textbf{割当更新}：
    $i$ を $j_i^1$ に割り当てる．
    既に $j_i^1$ が別の入札者 $i'$ に割り当てられていた場合は，
    $i'$ を未割当に戻す．
\end{enumerate}

全員が割り当てられたら終了する．

\begin{figure}[htbp]
\centering
\begin{tikzpicture}[>=Stealth, scale=0.85,
  person/.style={draw, circle, minimum size=18pt, fill=black!10},
  item/.style={draw, rectangle, minimum size=18pt, fill=black!25}]
  % People
  \node[person] (p1) at (0, 2) {\small $1$};
  \node[person] (p2) at (0, 0) {\small $2$};
  \node[person] (p3) at (0, -2) {\small $3$};
  \node[left=4pt] at (0, 3) {入札者};
  % Items
  \node[item] (i1) at (5, 2) {\small $1$};
  \node[item] (i2) at (5, 0) {\small $2$};
  \node[item] (i3) at (5, -2) {\small $3$};
  \node[right=4pt] at (5, 3) {品物};
  % Current assignment
  \draw[->, thick] (p1) -- (i2) node[midway, above, sloped]{\scriptsize 割当};
  \draw[->, thick, dashed] (p2) -- (i2)
    node[midway, below, sloped]{\scriptsize 入札};
  \draw[->, thick] (p3) -- (i3) node[midway, above, sloped]{\scriptsize 割当};
  % Price labels
  \node[right=10pt] at (i1.east) {$\mathbf{g}_1$};
  \node[right=10pt] at (i2.east) {$\mathbf{g}_2$};
  \node[right=10pt] at (i3.east) {$\mathbf{g}_3$};
\end{tikzpicture}
\caption{オークションアルゴリズムの1ステップ：
入札者 2 が品物 2 に入札し，現在の割当者（入札者 1）を追い出す．}
\label{fig:auction}
\end{figure}

\subsection{$\varepsilon$-相補性条件}\label{subsec:eps-cs}

パラメータ $\varepsilon > 0$ はアルゴリズムの精度と速度のトレードオフを制御する．
各反復で以下の\textbf{$\varepsilon$-相補性条件} ($\varepsilon$-CS) が保たれる：
\[
  \forall\, i \in S, \quad
  \mathbf{C}_{i,\xi_i} - \mathbf{g}_{\xi_i}
  \leq \varepsilon + \min_{j} (\mathbf{C}_{i,j} - \mathbf{g}_j).
\]
これは「各入札者が最良の品物から $\varepsilon$ 以内のコストの品物に
割り当てられている」ことを意味する．

\begin{proposition}{オークションアルゴリズムの停止性と精度}{auction-stop}
オークションアルゴリズムは高々
$N = n \|\mathbf{C}\|_\infty / \varepsilon$ ステップで停止し，
得られる割当のコストは最適値から $n\varepsilon$ 以内である．
\end{proposition}

$\varepsilon$\textbf{-スケーリング}と呼ばれる技法では，
大きな $\varepsilon$ から開始して
反復ごとに $\varepsilon$ を縮小していくことで，
計算量を $O(n^3 \log \|\mathbf{C}\|_\infty)$ に改善できる．


% ============================================================
\section{アルゴリズムの比較}\label{sec:algo-comparison}

本章で紹介した3つのアルゴリズムを比較する．

\medskip
\begin{center}
\begin{tabular}{lccc}
  & ネットワーク単体法 & 双対上昇法 & オークション \\
  \hline
  操作対象 & 主問題（$\mathbf{P}$） & 双対問題（$\mathbf{f}, \mathbf{g}$）
    & 双対（$\mathbf{g}$）+ 割当 \\
  解の種類 & 厳密解 & 厳密解 & $\varepsilon$-近似解 \\
  対象 & 一般の $\mathbf{a}, \mathbf{b}$ &
    一般の $\mathbf{a}, \mathbf{b}$ & 主に割当問題 \\
  計算量の目安 & $O(n^3 \log n)$ & $O(n^3)$ & $O(n^3 \log n)$ \\
  \hline
\end{tabular}
\end{center}
\medskip

いずれの手法も $n$ が大きい場合（数千点以上）では計算コストが大きく，
大規模な問題には次章で述べるエントロピー正則化に基づく
近似手法が必要となる．
